\documentclass[12pt,a4paper]{article}

\usepackage[utf8]{inputenc}
\usepackage[T1]{fontenc}
\usepackage[spanish,es-nodecimaldot]{babel}
\shorthandoff{"}
\usepackage{lmodern}
\usepackage{microtype}
\usepackage{geometry}
\geometry{top=2.5cm,bottom=2.5cm,left=3cm,right=2.5cm}
\usepackage{graphicx}
\usepackage{float}
\usepackage{caption}
\usepackage{booktabs}
\usepackage{array}
\usepackage{tabularx}
\usepackage{enumitem}
\usepackage{xcolor}
\usepackage{titlesec}
\usepackage{fancyhdr}
\usepackage{listings}
\usepackage{tikz}
\usetikzlibrary{positioning,arrows.meta,shapes.geometric,fit,backgrounds,calc,babel}
\usepackage[hidelinks,breaklinks]{hyperref}

\newcommand{\titulo}{Implementaci\'on de Jenkins como servidor de integraci\'on continua}
\newcommand{\materia}{\'Enfasis Profesional I (Integraci\'on Continua)}
\newcommand{\docente}{Edward Reyes Corredor}
\newcommand{\integrante}{Miguel Angel Aguilar Luna}
\newcommand{\usuario}{SavitarADSO}
\newcommand{\fecha}{24 de noviembre de 2026}
\newcommand{\repositorio}{\url{https://github.com/SavitarADSO/-nfasis-Profesional-I-Integraci-n-Continua-}}
\definecolor{negro}{HTML}{000000}
\definecolor{gris}{HTML}{595959}
\definecolor{fondo}{HTML}{F2F2F2}
\titleformat{\section}{\large\bfseries\color{negro}}{\thesection.}{0.6em}{}
\titleformat{\subsection}{\normalsize\bfseries\color{negro}}{\thesubsection}{0.6em}{}
\pagestyle{fancy}
\fancyhf{}
\setlength{\headheight}{15pt}
\renewcommand{\headrulewidth}{0.4pt}
\fancyhead[L]{\footnotesize Entrega 2 -- Jenkins}
\fancyhead[R]{\footnotesize \'Enfasis Profesional I}
\fancyfoot[C]{\small\thepage}
\setlength{\parindent}{0.6cm}
\setlength{\parskip}{0.35em}
\sloppy
\setlist[itemize]{leftmargin=1.5em,itemsep=0.2em,topsep=0.3em}
\setlist[enumerate]{leftmargin=1.8em,itemsep=0.2em,topsep=0.3em}
\lstset{
  basicstyle=\ttfamily\footnotesize,
  backgroundcolor=\color{fondo},
  numbers=left,
  numberstyle=\tiny\color{gris},
  numbersep=8pt,
  frame=single,
  rulecolor=\color{gris},
  breaklines=true,
  columns=fullflexible,
  keepspaces=true,
  showstringspaces=false
}
\newcommand{\captura}[2]{%
  \begin{figure}[H]
    \centering
    \IfFileExists{#1}{%
      \includegraphics[width=0.88\textwidth]{#1}%
    }{%
      \fbox{\parbox[c][4.2cm][c]{0.84\textwidth}{%
        \centering\color{gris}\textbf{\large EVIDENCIA PENDIENTE}\\[0.5em]
        \small Guarda la captura solicitada con este nombre:\\[0.4em]
        \texttt{#1}}}%
    }
    \caption{#2}
  \end{figure}
}

\begin{document}
\begin{titlepage}
\centering
\vspace*{1.4cm}
{\color{negro}\rule{\textwidth}{1.5pt}}\\[0.6cm]
{\LARGE\bfseries \titulo\par}
\vspace{0.6cm}
{\large Segunda entrega -- Proyecto de integraci\'on continua\par}
{\color{negro}\rule{\textwidth}{1.5pt}}
\vspace{1.8cm}
{\large\bfseries Integrante\par}
\vspace{0.3cm}
{\Large \integrante\par}
{\usuario\par}
\vspace{1.2cm}
{\large\bfseries Docente\par}
\vspace{0.3cm}
{\Large \docente\par}
\vfill
{\large \materia\par}
\vspace{0.4cm}
{\large \fecha\par}
\vspace{0.6cm}
{\small Repositorio: \repositorio\par}
\end{titlepage}

\tableofcontents
\newpage

\section{Introducci\'on}
En la primera entrega se implement\'o una API Flask que persiste mensajes en
PostgreSQL. Ambos servicios se orquestan con Docker Compose y la aplicaci\'on
encuentra la base de datos mediante el nombre de servicio \texttt{db}. Esta
segunda entrega incorpora Jenkins al mismo proyecto para automatizar la descarga
del c\'odigo, la construcci\'on de la imagen de la aplicaci\'on y una prueba de
sintaxis reproducible.

Jenkins se ejecuta localmente como un tercer servicio de Compose. El archivo
\texttt{Jenkinsfile}, versionado en la ra\'iz del repositorio, describe las
etapas del pipeline. Este informe documenta los requisitos, la arquitectura,
la configuraci\'on y las evidencias de la ejecuci\'on.

\section{Objetivos}
\subsection{Objetivo general}
Integrar Jenkins con el proyecto Docker existente y con GitHub para construir
y verificar la imagen de Flask mediante un pipeline como c\'odigo.

\subsection{Objetivos espec\'ificos}
\begin{itemize}
  \item Documentar los requisitos b\'asicos de hardware, software, red y almacenamiento.
  \item A\~nadir Jenkins al archivo Compose sin sustituir los servicios de la Entrega 1.
  \item Conservar de forma persistente la configuraci\'on de Jenkins.
  \item Configurar un job que obtenga el \texttt{Jenkinsfile} desde la rama principal.
  \item Ejecutar y registrar con evidencia las etapas Checkout, Build y Test.
\end{itemize}

\section{Requisitos para la implementaci\'on}
\subsection{Recursos y software}
Para una pr\'actica local conviene disponer de al menos 4 GB de memoria RAM
asignada a Docker Desktop y espacio libre para la imagen de Jenkins, las im\'agenes
del proyecto y el historial de ejecuciones. La primera construcci\'on descarga
la imagen base y los paquetes requeridos, por lo que puede tomar varios minutos.

\begin{table}[H]
\centering
\caption{Componentes y funci\'on en esta soluci\'on}
\begin{tabularx}{\textwidth}{>{\bfseries}p{3.2cm}X}
\toprule
Componente & Uso \\
\midrule
Docker Desktop y Compose v2 & Ejecutar y orquestar los tres servicios. \\
Jenkins 2.541.3 con JDK 17 & Servidor que administra el job y sus ejecuciones. \\
Docker CLI & Construir la imagen de Flask desde el pipeline. \\
GitHub p\'ublico & Fuente del c\'odigo y del archivo Jenkinsfile. \\
Navegador web & Configuraci\'on inicial de Jenkins en localhost:8080. \\
\bottomrule
\end{tabularx}
\end{table}

\subsection{Red, puertos y persistencia}
El servicio Jenkins comparte la red bridge \texttt{red\_proyecto} con \texttt{app}
y \texttt{db}. El puerto web 8080 se publica solo en la interfaz local del equipo;
50000 queda reservado para agentes entrantes. El volumen nombrado
\texttt{jenkins\_home} conserva usuarios, plugins, jobs y ejecuciones cuando el
contenedor se recrea.

\subsection{Acceso al motor Docker}
El pipeline necesita comunicarse con el motor Docker del anfitri\'on. Para esta
pr\'actica local se monta \texttt{/var/run/docker.sock} y Jenkins se ejecuta como
root dentro del contenedor. Este acceso equivale a conceder control sobre el
motor Docker; por ello la interfaz se enlaza a localhost y no se debe exponer
este montaje a un servidor o red no confiable.

\section{Arquitectura}
\begin{figure}[H]
\centering
\begin{tikzpicture}[
  node distance=1.5cm and 1.7cm,
  servicio/.style={draw,rounded corners,align=center,minimum width=3.0cm,minimum height=1.2cm},
  flecha/.style={-{Latex[length=2.5mm]},thick}
]
\node[servicio] (git) {GitHub\\Repositorio p\'ublico};
\node[servicio,right=of git] (jenkins) {Jenkins\\Checkout / Build / Test};
\node[servicio,below right=of jenkins] (app) {Flask\\Imagen construida};
\node[servicio,below left=of app] (db) {PostgreSQL\\Servicio db};
\node[draw,dashed,rounded corners,fit=(jenkins)(app)(db),inner sep=0.55cm,label=below:{Red Compose \texttt{red\_proyecto}}] (red) {};
\draw[flecha] (git) -- node[above,align=center]{Git HTTPS\\Poll SCM} (jenkins);
\draw[flecha] (jenkins) -- node[right,align=left]{Docker CLI\\socket del motor} (app);
\draw[flecha] (app) -- node[below,align=center]{DNS Compose: db:5432} (db);
\end{tikzpicture}
\caption{Vista general de la arquitectura de la Entrega 2.}
\end{figure}

GitHub aloja el c\'odigo y el pipeline. Jenkins consulta el repositorio y ejecuta
Docker CLI mediante el socket montado. Los servicios existentes Flask y PostgreSQL
mantienen su comunicaci\'on por el nombre \texttt{db}; Jenkins no cambia esa
configuraci\'on ni reemplaza GitHub Actions de la primera entrega.

\section{Implementaci\'on}
\subsection{Imagen y servicio Jenkins}
El Dockerfile utiliza una versi\'on identificada de Jenkins con JDK 17 e instala
Docker CLI. En Compose se agrega el servicio \texttt{jenkins}, su volumen
\texttt{jenkins\_home}, los puertos locales y el socket Docker. La aplicaci\'on y
la base de datos permanecen sin cambios.

\begin{lstlisting}[caption={Contenido de jenkins/Dockerfile}]
FROM jenkins/jenkins:2.541.3-lts-jdk17

USER root
RUN apt-get update \
    && apt-get install -y --no-install-recommends docker.io \
    && rm -rf /var/lib/apt/lists/*
USER jenkins
\end{lstlisting}

\begin{lstlisting}[caption={Extracto del servicio Jenkins en docker-compose.yml}]
jenkins:
  build: ./jenkins
  container_name: ci_jenkins
  user: root
  ports:
    - "127.0.0.1:8080:8080"
    - "127.0.0.1:50000:50000"
  volumes:
    - jenkins_home:/var/jenkins_home
    - /var/run/docker.sock:/var/run/docker.sock
  networks:
    - red_proyecto
\end{lstlisting}

\captura{figuras/entrega2/01-compose-jenkins.png}{Estado de los tres servicios en Docker Compose.}
\captura{figuras/entrega2/02-jenkins-logs.png}{Registro de inicio del servicio Jenkins.}

\subsection{Configuraci\'on inicial}
En el primer inicio, Jenkins solicita la contrase\~na temporal almacenada en
\texttt{/var/jenkins\_home/secrets/initialAdminPassword}. Despu\'es de desbloquear
el servidor se instalan los plugins sugeridos y se crea un usuario administrador.
La contrase\~na no forma parte del repositorio ni debe aparecer en las capturas.

\captura{figuras/entrega2/03-unlock-jenkins.png}{Asistente de desbloqueo con la contrase\~na oculta.}
\captura{figuras/entrega2/04-jenkins-inicializado.png}{Panel principal tras la configuraci\'on inicial.}

\subsection{Pipeline como c\'odigo}
El archivo \texttt{Jenkinsfile} define tres etapas. Checkout obtiene la revisi\'on
configurada en el job; Build crea la imagen de Flask; Test ejecuta el compilador
de Python dentro de esa imagen. Se deshabilitan ejecuciones simult\'aneas para
evitar que dos jobs reemplacen la misma etiqueta Docker a la vez.

\begin{lstlisting}[caption={Jenkinsfile}]
pipeline {
    agent any
    options {
        disableConcurrentBuilds()
        timestamps()
    }
    stages {
        stage('Checkout') {
            steps { checkout scm }
        }
        stage('Build') {
            steps {
                sh 'docker build -t proyecto-ci-app:jenkins ./app'
            }
        }
        stage('Test') {
            steps {
                sh 'docker run --rm --entrypoint python proyecto-ci-app:jenkins -m py_compile app.py'
            }
        }
    }
}
\end{lstlisting}

\captura{figuras/entrega2/05-jenkinsfile-github.png}{Archivo Jenkinsfile versionado en la ra\'iz del repositorio.}

\subsection{Job, SCM y sondeo}
En la interfaz de Jenkins se crea un elemento de tipo Pipeline. Se elige
\emph{Pipeline script from SCM}, SCM Git, la URL HTTPS del repositorio p\'ublico,
la rama \texttt{*/main} y la ruta \texttt{Jenkinsfile}. No se configuran
credenciales porque el repositorio es p\'ublico. Se activa Poll SCM con
\texttt{H/5 * * * *}; el sondeo permite detectar cambios aunque Jenkins se ejecute
en localhost, inaccesible para webhooks entrantes de GitHub.

\captura{figuras/entrega2/06-job-scm.png}{Job configurado para leer el pipeline desde GitHub.}
\captura{figuras/entrega2/07-poll-scm.png}{Disparador Poll SCM con intervalo aproximado de cinco minutos.}

\subsection{Ejecuci\'on}
Una ejecuci\'on correcta debe mostrar las etapas Checkout, Build y Test en verde.
La consola debe terminar en \texttt{Finished: SUCCESS}; ante un fallo, la consola
permite localizar si ocurri\'o durante la descarga, construcci\'on o prueba.

\captura{figuras/entrega2/08-pipeline-verde.png}{Vista de etapas con una ejecuci\'on completada correctamente.}
\captura{figuras/entrega2/09-console-success.png}{Final de la consola con el resultado SUCCESS.}

\section{Procedimiento de puesta en marcha}
\begin{lstlisting}[caption={Comandos locales}]
docker compose up -d --build
docker compose ps
docker compose logs -f jenkins
docker compose exec jenkins cat /var/jenkins_home/secrets/initialAdminPassword
\end{lstlisting}
Al terminar la configuraci\'on web, se guarda el job y se ejecuta \emph{Build
Now}. Las capturas se guardan en \texttt{docs/figuras/entrega2/} siguiendo la
lista de nombres de \texttt{LEEME.md}. Para detener los servicios sin perder los
datos de Jenkins se usa \texttt{docker compose down}; no se debe a\~nadir \texttt{-v}
si se quiere conservar su configuraci\'on.

\section{Conclusiones}
Jenkins se integra en la infraestructura existente como servicio adicional,
mientras Flask, PostgreSQL y el pipeline de GitHub Actions de la Entrega 1
conservan su funci\'on. El pipeline almacenado en el repositorio hace reproducible
la construcci\'on de la imagen y comprueba la sintaxis de la aplicaci\'on. El
volumen nombrado preserva la configuraci\'on del servidor; Poll SCM permite
consultar GitHub sin publicar Jenkins en Internet. El montaje del socket Docker
es una concesi\'on de privilegios significativa y queda limitado a la pr\'actica
local.

\section*{Referencias}
\begin{itemize}
  \item Jenkins. \emph{Jenkins User Documentation}. \url{https://www.jenkins.io/doc/}
  \item Jenkins. \emph{Installing Jenkins}. \url{https://www.jenkins.io/doc/book/installing/}
  \item Docker Inc. \emph{Docker Compose documentation}. \url{https://docs.docker.com/compose/}
  \item GitHub. \emph{About GitHub Actions}. \url{https://docs.github.com/en/actions}
\end{itemize}

\end{document}
